\hypertarget{ui-ux-spec-ascend-uiux-design-brief-for-the-web-client-pwa}{%
\subsection{Ascend: UI/UX design brief for the web client
(PWA)}\label{ui-ux-spec-ascend-uiux-design-brief-for-the-web-client-pwa}}

The original brief below is retained as design history. The
owner\textquotesingle s subsequent request for a more motivating
interface is implemented and documented in the
\href{engagement-audit.md}{engagement audit}, which supersedes the
logging-first flow and visual constraints here while retaining the
four-screen structure.

This is a precise design brief for whoever builds \texttt{web/} (see
\path|spec-v3-public-release.md| §10 Phase 2). Read it alongside that
file and \path|hackstreak-spec-v2.md|. It describes exactly what to
build and why, not just "make it look nice" --- follow the structure,
states and rules below rather than inventing your own IA.

\hypertarget{ui-ux-spec-0-the-one-thing-this-ui-must-not-become}{%
\subsubsection{0. The one thing this UI must not
become}\label{ui-ux-spec-0-the-one-thing-this-ui-must-not-become}}

The whole project exists because the owner gets overwhelmed by
open-ended choice and loses the thread. A UI that shows everything at
once --- all levels, all stats, all settings, on one dashboard ---
recreates the exact problem this tool is supposed to solve. \textbf{The
test for every screen: could someone open this app, in under 3 seconds,
see the one thing they\textquotesingle re meant to do right now, and do
it in one tap?} If a screen fails that test, it\textquotesingle s wrong,
no matter how complete or impressive it looks.

Second constraint, just as important: \textbf{don\textquotesingle t
build dark patterns.} No fake urgency ("only 2 hours left!"), no
guilt-laden copy ("You\textquotesingle re letting your streak die!"), no
manipulative streak-loss modals, no infinite-scroll feeds. The owner has
already noticed these mechanics pull on them harder than
they\textquotesingle d like. Design for a satisfying, honest "I did the
thing" feeling --- not for compulsive reopening.

\hypertarget{ui-ux-spec-1-information-architecture-4-screens-bottom-nav-nothing-deeper-than-2-taps}{%
\subsubsection{1. Information architecture: 4 screens, bottom nav,
nothing deeper than 2
taps}\label{ui-ux-spec-1-information-architecture-4-screens-bottom-nav-nothing-deeper-than-2-taps}}

\begin{verbatim}
[ Today ]   [ Path ]   [ Stats ]   [ More ]
\end{verbatim}

Four tabs, fixed bottom nav bar on mobile (thumb-reachable), persistent
left rail on desktop/tablet widths. Four, not five or six --- every
extra top-level destination is a decision the owner has to filter out
before reaching the thing they came for. Anything that
doesn\textquotesingle t earn a tab (settings, pack management, parked
ideas, account) lives inside \textbf{More}, not bolted onto the nav bar.

\begin{itemize}
\tightlist
\item
  \textbf{Today} --- the default screen on open. One primary action.
  This is 80\% of real usage.
\item
  \textbf{Path} --- the visual level map (the Duolingo-inspired piece),
  browsed, not acted from by default.
\item
  \textbf{Stats} --- streak, heatmap, minutes, search. Looked at, not
  acted on.
\item
  \textbf{More} --- settings, missions/packs, nudges, account.
  Everything infrequent.
\end{itemize}

\hypertarget{ui-ux-spec-2-today-screen-the-one-that-matters-most}{%
\subsubsection{2. Today screen (the one that matters
most)}\label{ui-ux-spec-2-today-screen-the-one-that-matters-most}}

Single-column, generous vertical spacing, nothing below the fold on a
phone except what\textquotesingle s described here.

\begin{enumerate}
\def\labelenumi{\arabic{enumi}.}
\tightlist
\item
  \textbf{Header strip} (small, top, never dominant): streak number + a
  small flame/leaf glyph, tappable → jumps to Stats. It must read as
  \emph{information}, not as a warning light. No red, no countdown, no
  "at risk" badge glowing here --- "at risk" framing is reserved for the
  nudge notification itself (spec v2 §11b), not for a screen the owner
  opens voluntarily. Being reminded gently when prompted is fine; being
  alarmed every time they check in is not.
\item
  \textbf{Time-available chips}, directly under the header:
  \texttt{10 min} \texttt{30 min} \texttt{2 hr+}
  \texttt{Surprise me}. Tapping one filters what "Next" shows below
  (implements \texttt{ascend \textless{}minutes\textgreater{}} from the
  spec) --- this is how a busy day becomes a zero-decision day, since
  the owner never has to mentally estimate what fits.
\item
  \textbf{The Next card} --- the hero of the screen, large, filled with
  the accent color, rounded corners, generous padding:

  \begin{itemize}
  \tightlist
  \item
    Mission name, small, muted, above the title (context, not the focus)
  \item
    Node title, large, bold
  \item
    One line: the \texttt{done\_when} sentence, in the
    owner\textquotesingle s actual words, not a category label
  \item
    A size pill (\texttt{\textasciitilde{}20 min}, derived from S/M/L,
    not the raw letter --- minutes mean something, a letter
    doesn\textquotesingle t)
  \item
    A milestone node gets a distinct visual treatment here too (a small
    star/diamond accent, a slightly different card border) so "this
    one\textquotesingle s the real test" is legible before they even tap
    in
  \item
    One primary button: \textbf{Log it}. Not "Start," not "Begin
    session" --- the action that matters is recording real work, and the
    button name should say exactly that.
  \end{itemize}
\item
  \textbf{The Quick win card}, directly below, visually secondary:
  smaller, outlined (not filled), lower visual weight. Same shape
  (title, one line, size pill, Log it button) but unmistakably "the
  alternate," not "the second mandatory task." This is the side-lane
  node (e.g. the next Bandit level). On a day the owner only has 10
  minutes and no energy for the main path, this is what they tap, with
  zero guilt implied by its smaller presence.
\item
  Nothing else on this screen. No activity feed, no "other packs you
  could start," no ads for features. If there\textquotesingle s a parked
  idea or a pending nudge-config issue, it does not appear here --- it
  waits in More.
\end{enumerate}

\hypertarget{ui-ux-spec-log-it-flow-bottom-sheet-not-a-new-page}{%
\paragraph{Log-it flow (bottom sheet, not a new
page)}\label{ui-ux-spec-log-it-flow-bottom-sheet-not-a-new-page}}

Tapping "Log it" slides up a sheet over the current screen (the owner
never loses their place):

\begin{itemize}
\tightlist
\item
  Two required fields only: \textbf{"What I did"} (free text, the
  \texttt{learned} field) and \textbf{"What\textquotesingle s next"}
  (\texttt{next\_step}), the latter pre-filled with a sensible default
  (the next node\textquotesingle s title) so it\textquotesingle s
  editable, not blank --- reduces the blank-page stall this owner has
  described elsewhere.
\item
  A collapsed \textbf{"Add details"} disclosure (closed by default)
  reveals: minutes spent, a proof link/file path, and --- only if the
  node has a verifier --- a "Mark as done" toggle with the
  verifier\textquotesingle s pass/fail shown inline.
\item
  One confirm button: \textbf{Log}. On success, the sheet closes with a
  brief, quiet confirmation (a checkmark that animates in and fades,
  under 400ms, not a confetti burst or a modal the owner must dismiss)
  and the streak number in the header ticks up. Small and satisfying,
  not a spectacle.
\end{itemize}

\hypertarget{ui-ux-spec-3-path-screen-the-duolingo-inspired-piece--adapted-not-copied}{%
\subsubsection{3. Path screen (the Duolingo-inspired piece --- adapted,
not
copied)}\label{ui-ux-spec-3-path-screen-the-duolingo-inspired-piece--adapted-not-copied}}

Duolingo\textquotesingle s winding horizontal path is built for one
fixed curriculum. This app has editable, user-authored levels of varying
length, so: \textbf{a simple vertical path, grouped by level, scrolling
top to bottom.} Easier to build, easier to scan, and it still delivers
the thing that actually matters from Duolingo\textquotesingle s design
--- visible structure and visible progress --- without copying its
visual identity.

\begin{itemize}
\tightlist
\item
  Levels render as sections in pack order, each with a header: level
  title, a thin progress bar, and \texttt{3/8 steps}. A level also
  shows its active-days count in small muted text (the day-wise tracking
  from the spec), e.g. \texttt{5 active days}.
\item
  Within a level, nodes sit on a vertical line as circles:

  \begin{itemize}
  \tightlist
  \item
    \textbf{Locked} --- grey outline, a small padlock glyph, muted title
    text. Tapping shows a brief tooltip: "Unlocks after {[}previous
    level{]}." Locked must look \emph{dormant}, not \emph{forbidden} ---
    no red, no X mark. The owner has described real shame around
    unfinished things; a locked node should read as "not yet," never as
    a failure.
  \item
    \textbf{Open} --- accent-colored outline, filled title text,
    tappable.
  \item
    \textbf{Done} --- filled solid in the accent color with a checkmark.
    Satisfying, final, no further affordance.
  \item
    \textbf{Milestone} --- a distinctly different shape, not just a
    bigger circle: a diamond or hexagon rather than a circle (reuse
    whatever your icon set already has --- the point is unmistakable
    silhouette, not elaborate art), in a second accent color used
    nowhere else in the UI. This node must be visually "the boss level"
    at a glance, because the spec gates the whole
    level\textquotesingle s progress on it --- the owner should
    immediately understand that this one node is different in kind, not
    just in size.
  \end{itemize}
\item
  The \textbf{side lane} (quick-win material like Bandit) is NOT mixed
  into this vertical list. It renders as its own small horizontal
  scrollable strip pinned above the main path, visually separated by a
  label ("Anytime") and a lighter background tint --- making clear
  it\textquotesingle s always available and never gates anything.
\item
  Tapping any open or done node opens a small detail sheet (same sheet
  component as Log-it): title, \texttt{done\_when}, resource link (opens
  externally), verify status if any, and Log it / Mark done / Skip
  actions.
\end{itemize}

\hypertarget{ui-ux-spec-4-stats-screen}{%
\subsubsection{4. Stats screen}\label{ui-ux-spec-4-stats-screen}}

Three stacked sections, each a card, no dashboard-grid of a dozen tiles:

\begin{enumerate}
\def\labelenumi{\arabic{enumi}.}
\tightlist
\item
  \textbf{Streak block}: current streak (large number), best streak
  (small, beside it), freeze availability as a single quiet line ("1
  freeze available this week"), not an icon-heavy meter.
\item
  \textbf{Heatmap}: a standard GitHub-style calendar heatmap, 12 weeks,
  in the accent color at varying opacity (not a second unrelated color
  scale) so it reads as "this app," not a borrowed component.
\item
  \textbf{Minutes}: a simple weekly bar chart of logged minutes (the
  field the earlier audit found was captured but never shown) --- four
  to eight bars, nothing fancier.
\end{enumerate}

A \textbf{search bar} sits above all three: "Search your notes." Results
are a plain reverse-chronological list of past \texttt{learned} lines,
each showing its date and which node it was logged against. This is the
one place the owner\textquotesingle s accumulated months of real notes
become useful again, so don\textquotesingle t bury it in More.

\hypertarget{ui-ux-spec-5-more-screen}{%
\subsubsection{5. More screen}\label{ui-ux-spec-5-more-screen}}

A plain, grouped settings list --- this screen is allowed to look like a
conventional settings page, because it\textquotesingle s visited rarely
and clarity beats cleverness here:

\begin{itemize}
\tightlist
\item
  \textbf{Active mission} (switch / view others)
\item
  \textbf{Packs} (installed list, install new, the lint/review state of
  each)
\item
  \textbf{Reminders} (nudge time, channel: ntfy/Telegram/desktop,
  per-mission override, a "send test" button)
\item
  \textbf{Appearance} (light / dark / system --- see §6)
\item
  \textbf{Parked ideas} (the rabbit-hole list --- reviewed here,
  deliberately out of the way of Today)
\item
  \textbf{Account / Server} (only shown in server mode: who
  you\textquotesingle re logged in as, which server)
\end{itemize}

\hypertarget{ui-ux-spec-6-visual-language-concrete-values-not-vibes}{%
\subsubsection{6. Visual language (concrete values, not
vibes)}\label{ui-ux-spec-6-visual-language-concrete-values-not-vibes}}

\begin{itemize}
\tightlist
\item
  \textbf{Palette}: one calm accent, not a primary-color riot. A
  desaturated teal or forest green (\texttt{\#2F6F5E}-ish range) reads
  as "growth," fits the name "Ascend," and is calmer than
  Duolingo\textquotesingle s saturated greens/golds. Pair it with a
  near-white (\texttt{\#FAFAF8}) background in light mode and a
  near-black, not pure-black (\texttt{\#15181B}-ish) in dark mode ---
  pure black against bright accent colors reads harsher than it needs
  to. The milestone accent (used ONLY for milestone nodes and nowhere
  else) should be a warm, muted amber, not neon gold --- special, not
  loud.
\item
  \textbf{Type}: one typeface, a standard humanist sans (system font
  stack is fine --- this is a utility app, not a brand showcase). Clear
  hierarchy: large/bold for the Next card title and the streak number,
  everything else restrained. Don\textquotesingle t use more than 3 font
  sizes on any single screen.
\item
  \textbf{Spacing}: generous, not cramped --- this is an anti-clutter
  requirement, not a cosmetic one. Minimum 16px side gutters on mobile,
  at least 12px between stacked cards, comfortable 44px+ tap targets
  throughout (this also happens to satisfy basic accessibility
  tap-target guidance for free --- do this by default, it costs nothing,
  even though a full accessibility pass is out of scope per the earlier
  audit).
\item
  \textbf{Motion}: small and purposeful only --- the Log-it confirmation
  checkmark, a progress bar filling, a node flipping from open to done.
  Nothing animates to grab attention; nothing loops; no celebratory
  full-screen takeovers.
\item
  \textbf{Corners/elevation}: soft rounded corners (12-16px) on cards,
  one subtle shadow level, no heavy skeuomorphism and no
  flat-and-bordered brutalism either --- aim for the restrained-card
  look of a well-made note-taking app (closer to
  Trilium\textquotesingle s calm utility than a marketing landing page).
\end{itemize}

\hypertarget{ui-ux-spec-7-responsive-behavior-three-real-breakpoints-not-mobile-app-that-also-loads-on-desktop}{%
\subsubsection{7. Responsive behavior: three real breakpoints, not
"mobile app that also loads on
desktop"}\label{ui-ux-spec-7-responsive-behavior-three-real-breakpoints-not-mobile-app-that-also-loads-on-desktop}}

This is opened from a phone home screen some days and a laptop browser
other days --- design both as genuine primary use, not one real target
and one that merely "doesn\textquotesingle t break." Same 4-screen IA
throughout; what changes is layout density and nav placement, not the
structure.

\textbf{Mobile (\textless{} 640px)}

\begin{itemize}
\tightlist
\item
  Bottom nav bar (§1), single column, full-width cards, Log-it as a
  bottom sheet.
\item
  Side gutters 16px minimum (§6).
\end{itemize}

\textbf{Tablet (640--1024px)}

\begin{itemize}
\tightlist
\item
  Left nav rail replaces the bottom bar (icon + label each).
\item
  Today: Next and Quick-win cards sit side by side, not stacked ---
  there\textquotesingle s room for both without crowding.
\item
  Log-it opens as a centered modal instead of a bottom sheet.
\end{itemize}

\textbf{Desktop (1024px+)}

\begin{itemize}
\tightlist
\item
  Left nav rail, same as tablet.
\item
  Content area has a max width (roughly 920--1000px) and stays centered.
  Do not stretch cards edge-to-edge on a wide monitor --- extra width
  becomes margin, not bigger cards. A single card stretched wall-to-wall
  on a 27" monitor reads worse, not better.
\item
  Today: Next and Quick-win side by side, time-available chips above
  both.
\item
  Path: keep the vertical level path in the same centered max-width
  column as everything else --- do not turn it into a multi-column grid
  of levels. A taller, narrower scroll reads better than a wide, shallow
  one for a sequential roadmap; resist the obvious "I have width, let me
  use a grid" instinct here specifically.
\item
  Add hover states on interactive elements (cards, nodes, buttons) --- a
  subtle elevation or border-color shift is enough, nothing elaborate;
  desktop has a pointer, so static elements shouldn\textquotesingle t
  feel inert.
\item
  Worth including, nearly free: a keyboard shortcut (e.g. \texttt{L})
  that opens Log-it for the current Next step directly, for a user
  already at a keyboard. Optional, not load-bearing.
\end{itemize}

Standard CSS (flex/grid + media queries, or a utility
framework\textquotesingle s responsive classes) covers all of this ---
no JavaScript-driven layout switching needed.

\hypertarget{ui-ux-spec-8-what-this-brief-deliberately-leaves-out}{%
\subsubsection{8. What this brief deliberately leaves
out}\label{ui-ux-spec-8-what-this-brief-deliberately-leaves-out}}

No onboarding-flow screens, no animated empty-state illustrations, no
custom icon set (use a standard open icon library --- Lucide or Feather
are both fine, pick one and use it consistently), no marketing/landing
page. Those are real design work but not what makes this tool work day
to day, and building them now would be the same mistake as over-building
the backend before Phase 1 proved itself --- effort spent on polish
before the core loop is in daily use.

\hypertarget{ui-ux-spec-9-text-to-hand-the-coding-agent}{%
\subsubsection{9. Text to hand the coding
agent}\label{ui-ux-spec-9-text-to-hand-the-coding-agent}}

\begin{quote}
Build \texttt{web/} as described in \texttt{ui-ux-spec.md} precisely ---
the screen structure (Today/Path/Stats/More), the Next and Quick-win
card distinction, the vertical level path with
locked/open/done/milestone node states, the Log-it bottom sheet, and the
color/spacing/motion values given are not suggestions to riff on, they
are the spec. Where something is genuinely unspecified (icon choice,
exact shadow values), use your own judgment within the stated
constraints (calm palette, generous spacing, restrained motion) rather
than defaulting to a generic admin-dashboard look. Do not add screens,
nav items, or dashboard widgets beyond what\textquotesingle s listed.
Ask before deviating from the information architecture in section 1.
\end{quote}
